\SourceLeaf{194}mit diesem Gegenstand beschäftigt haben. Bei mir sind es über 30 Jahre, und ich glaube nicht,
dass sich jemand gerade mit diesem $2^{\mathrm n}.$ Teil mehr beschäftigt haben kann als ich,
obwohl ich niemals etwas darüber veröffentlicht habe. Die Annahme, dass die Summe der 3
Winkel kleiner als $180^{0}$ ist, führt zu einer eigenen, von unserer (euklidischen) ganz
verschiedenen Geometrie, die in sich durchweg konsequent ist und die ich für mich
selbst ganz befriedigend entwickelt habe, sodass ich jede Aufgabe darin lösen kann, mit
Ausnahme der Bestimmung einer Konstanten, die sich a priori nicht ermitteln lässt. Je größer
man diese Konstante annimmt, desto näher kommt man der euklidischen Geometrie, und ein
unendlich großer Wert lässt beide zusammenfallen. Die Sätze jener Geometrie scheinen zum
Teil paradox und dem Ungeübten ungereimt. Bei genauerer, ruhiger Überlegung findet man aber,
dass sie an sich durchaus nichts Unmögliches enthalten. So können zum Beispiel die drei Winkel eines
Dreiecks beliebig klein werden, wenn man nur die Seiten groß genug wählen darf.
Dennoch kann der Flächeninhalt eines Dreiecks, wie groß man die Seiten auch wählt, nie
eine bestimmte Grenze überschreiten, ja sie nicht einmal erreichen. Alle meine Bemühungen,
einen Widerspruch oder eine Inkonsequenz in dieser nichteuklidischen Geometrie zu finden, sind
fruchtlos geblieben. Das Einzige, was unserem Verstand daran widerstrebt, ist, dass es, wäre
sie wahr, im Raum eine \textls{an sich bestimmte} (wenn auch uns unbekannte) lineare Größe geben
müsste. Aber mir scheint, trotz der nichtssagenden Wortweisheit der Metaphysiker wissen wir
eigentlich zu wenig oder gar nichts über das wahre Wesen des Raums, als dass wir etwas uns
unnatürlich Vorkommendes mit dem \textls{absolut Unmöglichen} verwechseln dürften. Wäre die
nichteuklidische Geometrie die wahre und stünde jene Konstante in irgendeinem Verhältnis zu solchen
Größen, die im Bereich unserer Messungen auf der Erde oder am Himmel liegen, so ließe sie
sich a posteriori ermitteln. Ich habe daher gelegentlich im Scherz den Wunsch geäußert,
dass die euklidische Geometrie nicht die wahre wäre, weil wir dann ein absolutes Maß a priori
hätten.

Bei einem Mann, der sich mir als ein denkender mathematischer Kopf gezeigt hat, fürchte ich
nicht, dass er das Vorstehende missversteht. Auf jeden Fall sollen Sie es nur als
eine private Mitteilung ansehen. Sie dürfen sie weder öffentlich verwenden noch in einer Weise, die sie an die
Öffentlichkeit\space%
